\chapter{Il fattore temporale \texorpdfstring{$f_i(t)$}{fi(t)}}
\label{ch:fattore}

Il fattore $f_i(t)$ modula la penalità di fascia oraria (termine~2 del
capitolo~\ref{ch:obiettivo}). È l'elemento del modello che traduce nel modo più diretto le
dichiarazioni orarie del questionario, e ha due componenti indipendenti.

\section{Formulazione completa}

\begin{modelbox}{Fattore temporale}
Sia $d_i(t)$ la distanza \textbf{circolare}, in minuti, fra lo slot $t$ e la fascia dichiarata più
vicina dell'apparecchio $i$ (nulla se $t$ cade dentro la fascia). Allora
\begin{equation}
  f_i(t) =
  \begin{cases}
    w \cdot \vag_i
    & \text{se } t \text{ è dentro la fascia dichiarata,}\\[6pt]
    w \cdot \vag_i \;+\; \min\!\left( F_{\max},\; 1 + \dfrac{d_i(t)}{\rho} \right)
    & \text{altrimenti,}
  \end{cases}
  \label{eq:fattore}
\end{equation}
dove $\rho$ è \code{window\_penalty\_ramp\_min} (60 minuti), $F_{\max}$ è
\code{window\_penalty\_max\_factor} (6) e $w$ è \code{window\_vagueness\_weight}.
\end{modelbox}

\section{Prima componente: distanza dalla fascia}

Il secondo addendo della~\eqref{eq:fattore} penalizza lo scostamento dall'orario dichiarato, in
modo \textbf{graduale}.

\paragraph{Perché graduale.} Una penalità piatta tratterebbe allo stesso modo un'attivazione dieci
minuti dopo la fine della fascia e una a otto ore di distanza. La prima è del tutto plausibile ---
l'utente ha indicato un orario di massima, non un vincolo --- la seconda contraddice la
dichiarazione. La formulazione a rampa distingue i due casi.

\paragraph{Comportamento.} Sul bordo della fascia il fattore vale 1, cioè la penalità base; cresce
di un'unità ogni $\rho = 60$ minuti di allontanamento; satura a $F_{\max} = 6$ anziché divergere.
La saturazione evita che apparecchi molto distanti dalla propria fascia ricevano penalità
arbitrariamente grandi, che li renderebbero di fatto inutilizzabili e comprometterebbero la
ricostruzione nelle ore scoperte.

\paragraph{Distanza circolare.} La distanza è calcolata sul quadrante delle 24 ore, non sulla retta
dei minuti. Questo è necessario perché molte fasce dichiarate scavalcano la mezzanotte: per una
fascia $23{:}00$--$02{:}00$, le $01{:}00$ sono \emph{dentro} e le $22{:}30$ distano trenta minuti,
non ventitré ore e mezza.

\section{Seconda componente: genericità della dichiarazione}

Il primo addendo della~\eqref{eq:fattore} è un \textbf{pavimento} costante, presente anche dentro
la fascia dichiarata.

\subsection{Il problema che risolve}

Nella formulazione senza pavimento ($w = 0$) si verifica un'asimmetria non voluta:

\begin{center}
\begin{tabularx}{\textwidth}{@{}l l X@{}}
\toprule
\textbf{Apparecchio} & \textbf{Costo} & \textbf{Situazione} \\
\midrule
con fascia, dentro    & $0$        & non paga nulla \\
con fascia, fuori     & $\ge \lambda_w p_i$ & paga la penalità graduata \\
\textbf{senza fascia} & $\mathbf{0}$ & \textbf{non paga nulla, a nessuna ora} \\
\bottomrule
\end{tabularx}
\end{center}

Un apparecchio che non ha dichiarato alcun orario è quindi un \textbf{jolly gratuito}: può
assorbire qualsiasi plateau a qualsiasi ora senza costo. Nella fascia serale dichiarata dalla
lavastoviglie, lavastoviglie e climatizzatore pagano entrambi zero: pareggiano, e la scelta fra
loro è arbitraria --- decisa da dettagli numerici privi di significato fisico.

\subsection{L'impostazione corretta}

La soluzione \textbf{non} consiste nel penalizzare l'assenza di dichiarazione. Molti apparecchi
non hanno realmente un orario tipico --- climatizzatore, televisore, computer si usano a qualsiasi
ora --- e l'assenza di dichiarazione è \emph{assenza di informazione}, non evidenza di
implausibilità. Penalizzarli in quanto tali sarebbe scorretto.

L'impostazione adottata è diversa: \textbf{``nessuna fascia'' viene letta come ``fascia di 24
ore''}, che è letteralmente ciò che significa, e si penalizza quanto poco la dichiarazione
restringe il campo.

\begin{modelbox}{Misura di genericità}
\begin{equation}
  \vag_i \;=\; \frac{\log W_i}{\log 24} \;\in\; [0,1],
  \label{eq:vaghezza}
\end{equation}
dove $W_i$ è l'ampiezza complessiva in ore delle fasce dichiarate, posta a 24 in assenza di
dichiarazione.
\end{modelbox}

\subsection{Da dove viene la forma logaritmica}

La~\eqref{eq:vaghezza} non è una scelta arbitraria ma discende dalla verosimiglianza. Si modelli
l'orario di avvio dell'apparecchio come una variabile aleatoria uniforme sulla fascia dichiarata:
la densità di probabilità è $1/W_i$. Il costo informativo di affermare che l'apparecchio è attivo
in un dato istante è allora
\begin{equation}
  -\log \frac{1}{W_i} \;=\; \log W_i ,
\end{equation}
e la normalizzazione per $\log 24$ riporta la scala all'intervallo $[0,1]$, con una fascia di
un'ora a 0 e l'assenza di dichiarazione a 1.

L'interpretazione è la seguente: \textbf{chi si sbilancia su poche ore formula un'affermazione più
falsificabile, e riceve uno sconto quando l'evento cade effettivamente dentro la fascia
dichiarata}. Chi non si sbilancia non viene punito, semplicemente non riceve alcuno sconto.

\subsection{Valori sul campione reale}

\begin{center}
\begin{tabular}{@{}l l r r@{}}
\toprule
\textbf{Apparecchio} & \textbf{Fascia dichiarata} & $W_i$ & $\vag_i$ \\
\midrule
Lavastoviglie   & $19{:}00$--$22{:}00$ & 3 h  & 0.35 \\
Lavatrice       & $17{:}00$--$22{:}00$ & 5 h  & 0.51 \\
Lavatrice       & $10{:}00$--$17{:}00$ & 7 h  & 0.61 \\
Climatizzatore  & nessuna              & 24 h & 1.00 \\
Asciugatrice    & nessuna              & 24 h & 1.00 \\
\bottomrule
\end{tabular}
\end{center}

Si osservi che una fascia di 7 ore ottiene $0.61$, non molto lontano da $1.00$: la scala
logaritmica riconosce che una dichiarazione così ampia è \emph{quasi} priva di informazione, il che
è corretto.

\subsection{Effetto sull'ordinamento}

La tabella seguente riporta il fattore $f_i(t)$ con $w = 1$ per una lavastoviglie con fascia
$19{:}00$--$22{:}00$ e un climatizzatore senza fascia dichiarata.

\begin{center}
\begin{tabular}{@{}l r r l@{}}
\toprule
\textbf{Ora} & \textbf{Lavastoviglie} & \textbf{Climatizzatore} & \textbf{Preferito} \\
\midrule
$10{:}00$ & 6.35 & 1.00 & climatizzatore \\
$19{:}00$ & 0.35 & 1.00 & \textbf{lavastoviglie} \\
$21{:}00$ & 0.35 & 1.00 & \textbf{lavastoviglie} \\
$22{:}00$ & 1.36 & 1.00 & climatizzatore \\
\bottomrule
\end{tabular}
\end{center}

Il comportamento è quello desiderato in entrambe le direzioni: l'apparecchio specifico prevale
\textbf{soltanto dentro la fascia che ha dichiarato}, e l'ordinamento si inverte correttamente
appena se ne esce. Il pavimento non rende la lavastoviglie preferibile ovunque, ma solo dove la
sua dichiarazione lo giustifica.

\section{Retrocompatibilità}

Con $w = 0$ il pavimento scompare e la~\eqref{eq:fattore} si riduce alla formulazione originale:
fattore nullo dentro la fascia e penalità graduata soltanto fuori, con gli apparecchi senza fascia
che non pagano nulla. Tutti i modelli dal paragrafo~\ref{sec:m1} al~\ref{sec:m5} usano questa
impostazione; solo \code{cvxpy\_survey\_prior} attiva il pavimento.
